\section{Agent Externalization: consolidar las capacidades fuera de los parámetros}

El informe de Liu Weiwen (刘卫文) aborda la ruta que va de model scaling a agent evolution: muchas capacidades no necesitan comprimirse por completo en los parámetros del modelo base, sino que pueden externalizarse como memory, tool, skill, workflow y un harness que se puede modificar de manera sostenida. La externalization reduce el costo de actualización y también facilita auditar y combinar las capacidades.

\sourcefigure{0.93}{lecture07/slides-images/slide-002-00015s.jpg}{La próxima generación de Agent: personalización, proactividad y autoevolución}{00:00:15--00:01:00}

\sourcefigure{0.93}{lecture07/slides-images/slide-003-00060s.jpg}{Postura: agent evolution necesita nuevos portadores de capacidad}{00:01:00--00:04:00}

\subsection{De Model Scaling a Agent Evolution}

El modelo base aporta un prior general, pero las mejoras posteriores al despliegue suelen provenir de componentes externos: la memoria de recuperación complementa hechos nuevos, las herramientas aportan capacidades ejecutables, la biblioteca de skills conserva los flujos exitosos y el harness determina el orden de invocación. La policy efectiva del sistema es, en consecuencia, una función conjunta del modelo y del estado externo.

\sourcefigure{0.93}{lecture07/slides-images/slide-004-00240s.jpg}{De la expansión de parámetros a la coevolución modelo—Agent}{00:04:00--00:06:30}

\sourcefigure{0.93}{lecture07/slides-images/slide-005-00390s.jpg}{Ruta de investigación de la coevolución model--harness}{00:06:30--00:09:00}

\begin{importantbox}{Criterio para decidir la externalization}
El conocimiento que necesita actualizarse con frecuencia, que debe ser interpretable, o que se puede eliminar o combinar, es más adecuado para externalizarse; los patrones estables y ampliamente reutilizables son más adecuados para entrar a los parámetros. Un límite mal trazado provoca sobrecarga de contexto o bien obliga a reentrenar el modelo grande ante cada cambio pequeño.
\end{importantbox}

\subsection{BALTO: de texto de reflexión a conducta reutilizable}

Los métodos del tipo BALTO hacen que el Agent produzca una reflection a partir de trayectorias fallidas, y luego usa esa reflexión en la siguiente ronda de acción. Lo esencial no es generar frases vacías del tipo «la próxima vez ten más cuidado», sino extraer reglas que se puedan activar, ejecutar y verificar, y evaluar si transfieren entre tareas.

\sourcefigure{0.93}{lecture07/slides-images/slide-006-00540s.jpg}{Ciclo cerrado de self-evolving reflection de BALTO}{00:09:00--00:12:00}

\sourcefigure{0.93}{lecture07/slides-images/slide-007-00720s.jpg}{Balanced Rehearsal evita que la experiencia nueva sobrescriba capacidades previas}{00:12:00--00:13:30}

\begin{warningbox}{Memory pollution}
Si cada fallo se escribe en la memoria de largo plazo, la atribución errónea de causas, las experiencias fortuitas y las reglas que se contradicen entre sí contaminan rápidamente el context. El sistema de memoria necesita mecanismos de confianza, alcance, versión, deduplicación y descarte, y debe verificar lo que escribe mediante ejecución real.
\end{warningbox}

\subsection{LongTool: hacer que la capacidad de herramientas pueda crecer}

Aprender a usar herramientas no consiste solo en elegir entre las API existentes; también incluye comprender documentación extensa de herramientas, combinar varias herramientas y consolidar las invocaciones exitosas como un reusable skill. LongTool se enfoca en el aprendizaje y la generalización en entornos de herramientas de contexto largo, y muestra que el propio ecosistema externo de herramientas puede convertirse en una biblioteca de capacidades para la evolución del Agent.

\sourcefigure{0.93}{lecture07/slides-images/slide-009-00840s.jpg}{Marco de autoevolución del modelo de LongTool}{00:14:00--00:14:30}

\sourcefigure{0.93}{lecture07/slides-images/slide-010-00870s.jpg}{Resultados experimentales de LongTool y generalización de herramientas}{00:14:30--00:16:15}

\subsection{SKiMAS y MIMSkill: el harness y las skills multimodales también se pueden aprender}

SKiMAS trata el harness como un objeto modificable y hace que el sistema actualice la orquestación de módulos a partir de la interacción; MIMSkill, por su parte, transfiere skills de texto a tareas multimodales. Ambos muestran que un skill puede convertirse en una representación intermedia que conecta la policy abstracta con el entorno concreto.

\sourcefigure{0.93}{lecture07/slides-images/slide-011-00975s.jpg}{SKiMAS: hacer que el harness del Agent evolucione con la experiencia}{00:16:15--00:18:30}

\sourcefigure{0.93}{lecture07/slides-images/slide-012-01110s.jpg}{Marco modular de SKiMAS y su ruta de actualización}{00:18:30--00:19:30}

\sourcefigure{0.93}{lecture07/slides-images/slide-014-01365s.jpg}{MIMSkill: transferencia de skills de texto a skills multimodales}{00:22:45--00:25:15}

\begin{knowledgebox}{Definición técnica de un skill}
Un skill reutilizable debe incluir al menos una condición de activación, las entradas necesarias, una secuencia de acciones, una recuperación ante fallos y un criterio de terminación. Una «skill» que solo consiste en un resumen en lenguaje natural, sin una interfaz de ejecución, difícilmente logra integrarse de manera estable en la planificación automática.
\end{knowledgebox}

\subsection{ProAct: de la respuesta pasiva a la interacción proactiva}

Cuando el Agent puede mantener objetivos de largo plazo y el estado del entorno, ya no necesita esperar a que el usuario formule cada pregunta de manera explícita, y puede recordar algo, recopilar información o ejecutar acciones de preparación en el momento adecuado. ProAct estudia la proactive interaction, pero al mismo tiempo expone nuevos problemas de gobernanza: la conducta proactiva debe controlar la interrupción, los permisos y el riesgo irreversible.

\sourcefigure{0.93}{lecture07/slides-images/slide-016-01560s.jpg}{Marco de interacción proactiva de ProAct}{00:26:00--00:28:30}

\sourcefigure{0.93}{lecture07/slides-images/slide-018-01845s.jpg}{Experimento de interacción proactiva: cuándo actuar y cómo obtener retroalimentación del usuario}{00:30:45--00:32:15}

\subsection{Resumen de la sección}

La externalization de capacidades permite que el Agent aprenda de manera continua con un costo menor: la memory conserva la experiencia, la tool amplía la acción, el skill consolida los flujos, el harness organiza los componentes y la proactive policy decide cuándo actuar de manera proactiva. Un self-evolving system genuino debe estudiar a la vez la model evolution y la harness evolution, y dotar al estado externo de mecanismos de gobernanza y reversión.
